% Frozen before running the holdout grid below. Compile with xelatex.
\documentclass[11pt,a4paper]{article}
\usepackage{kotex}
\usepackage{amsmath,amssymb,booktabs,geometry}
\geometry{margin=25mm}
\title{I01 선검증 프로토콜: 관측 rank 개입의 적용 범위}
\author{RELATIONAL ATLAS}
\date{2026년 10월 4일}
\begin{document}
\maketitle

\section{목적과 보류할 주장}
본 검증은 ``유한한 위상 기준과 시계 아래에서 관측 차원의 식별
가능성''을 일반적인 결과로 주장하기 \emph{전}에 실시한다.
이미 관찰한 $n=4$, 9표지, 특정 상태 가중치의 rank $4/3$
패턴을 새로운 조건에 그대로 전이할 수 있는지, 어디에서 실패하는지
사전에 고정한 격자로 검사한다. 세 공간 차원의 자발적 선택은
검증 대상이나 결론이 아니다.

\section{고정한 홀드아웃 격자}
운동량 cutoff 1, $w_k\propto\exp(-\beta\sum_i k_i^2)$,
검출 축 $i$의 바이어스 $b$, 공통 위상 표지
$\phi_t=2\pi t/L$, 시계 에너지 $E(k)=\sum_i k_i^2$를 사용한다.
사전에 정한 조합은
\begin{equation}
 n\in\{3,4,5,6\},\quad
 \beta\in\{0.45,1.20\},\quad
 b\in\{0.55,1.25\},\quad
 \tau/\pi\in\{0,1/2\},\quad
 L\in\{2n-1,2n+1\}.
\end{equation}
각 조합에서 비제약 상태와 준비된 $P=0$ 상태의 표지 보존/삭제
고전 Fisher 행렬을 계산한다. 표지 삭제는 동일한 표본의
$t$ 정보를 합치는 분석이며 추가 물리 제약이 아니다.

\section{사전 기대와 판정 기준}
충분한 위상 격자 $L=2n+1$에서는 이 cutoff의 가능한 전하차
$|P-P'|\le2n$을 모두 평균하므로, 비제약 상태의 표지 삭제는
공통 위상을 정확히 없애야 한다. 사전 기대는 다음과 같다.
\begin{center}
\begin{tabular}{@{}lcc@{}}
\toprule
상태 & 표지 보존 & 표지 삭제, $L=2n+1$ \\
\midrule
비제약 & $n$ & $n-1$ \\
$P=0$ & $n-1$ & $n-1$ \\
\bottomrule
\end{tabular}
\end{center}
짧은 격자 $L=2n-1$에는 전하차 에일리어싱이 가능하다.
따라서 삭제 후 rank $n-1$이 깨지는지 숨기지 않고 경계 조건으로
기록한다. 판정은 공통방향 정보
$\boldsymbol1^T F\boldsymbol1/n$의 절댓값과 양의 고유값의
상대 크기를 함께 보고, 단일 고정 임계값만으로 rank를 주장하지
않는다. 수치적 정확 영은 별도 분석으로 검증한다.

\section{독립 검증 경로와 실패 처리}
모든 격자점의 해석적 확률/미분을 독립 중앙차분으로 교차 확인하고,
선택된 $n=3,4$ 상태는 Qiskit 회로의 직접 상태벡터 및 별도 Aer
계산 방법으로 확률을 비교한다. $P=0$ 공통 위상의 정확 불변성과
Aer 유한 샷 검정은 다른 질문으로 분리한다. 앞서 32개 명목
5\% 대조검정 중 5개가 양성이었던 현상은 독립 seed를 추가하고
정확 이항 null과 비교하되, 작은 $p$값을 물리 신호라고 해석하지
않는다. 불일치가 나오면 모형·회로·분석의 어느 경로에서 생겼는지
먼저 추적하고, 격자를 결과에 맞춰 바꾸지 않는다.

이 선검증이 끝나기 전에는 일반화된 ``관측 차원의 식별 가능성''을
확립된 결과나 논문 제목으로 사용하지 않는다. QPU 실험도 보류한다.
\end{document}
